%----------------------------------------------------------------------------------------------------
% START OF CHAPTER
%----------------------------------------------------------------------------------------------------
\chapter{Jan 31}
\paragraph{Topics} 
\textit{ATPG. Pronunciation. Possessive adjectives. Sentences.}

\section{ATPG: All the Possibilities Game}

\begin{table}[ht]
\centering
    \begin{tabular}{ r l  | r l  | r l }
         \itb{tavolo} & \iti{table}            & \itb{prenotatzione} & \iti{reservation}              & \itb{gelato} & \iti{gelato}  \\
         \hline
         the          & \iti{il tavolo}        & the                 & \iti{la prenotatzione}         & the          & \iti{il gelato} \\
         a            & \iti{un tavolo}        & a                   & \iti{una prenotatzione}        & a            & \iti{un gelato} \\
         two          & \iti{due tavoli}       & two                 & \iti{due prenotatzione}        & two          & \iti{due gelati} \\
         good         & \iti{un buon tavolo}   & good                & \iti{una buona prenotatzione}  & good         & \iti{un buon gelato} \\
         expensive    & \iti{un tavolo caro}   & expensive           & \iti{una prenotatzione cara}   & expensive    & \iti{un gelato caro} \\         
         new          & \iti{un nuovo tavolo}  & new                 & \iti{una nuova prenotatzione}  & new          & \iti{un nuovo gelato} \\
         big          & \iti{un tavolo grande} & big                 & \iti{una prenotatizone grande} & big          & \iti{un gelato grande} \\
         my           & \iti{il mio tavolo}    & my                  & \iti{la mia prenotatzione}     & my           & \iti{il mio gelato} \\
         your         & \iti{il tuo tavolo}    & your                & \iti{la tua prenotatzione}     & your         & \iti{il tuo gelato} \\
         his/her      & \iti{il suo tavolo}    & his/her             & \iti{la sua prenotatzione}     & his/her      & \iti{il suo gelato}
    \end{tabular}
    \caption{ATPG}
    \label{tab:tATPG-Jan31}
\end{table}
\section{Pronunciation}
Most consonants are pronounced as in English, but there are some rules to note:
\begin{enumerate}
    \item M / N / P / Q / R / S / T are identical.
    \item H at the start of a word is silent.
    \item Vowels \iti{a}, \iti{e}, and \iti{i} are pronounced like English `ah', `eh', and `ih'.
    \item Pairs of the same consonants are emphasized (which implicitly induces a tiny pause 
	    when spoken): \iti{bello} is pronounced like \iti{bel-lo}, not \iti{belo}.
    \item Letters \iti{c-} and \iti{g-} normally have a hard sound with their following vowel, 
	    e.g. \iti{ca}, \iti{go}, \iti{gu} sound as they do in English.
    \item However, the \iti{ce} and \iti{ci} pairs are `ch'-sounds, like English `cello' or `cheese'.
    \item Similarly, the \iti{ge} and \iti{gi} pairs are `j'-sounds, like English `general'.
\end{enumerate}
\begin{mdframed}
    \textbf{Watch out} for \iti{ch} and \iti{gh} combinations: the presence of the \iti{h} makes 
	for a hard sound. The food \iti{bruschetta} is \textit{broos-KET-ta} not 
	\textit{broo-SHET-uh}, no matter what the server at Olive Garden says. Another example: 
	\iti{pesce}, \textit{PEH-sheh} is `fish' but \iti{pesche}, \textit{PES-keh} is `peaches'.
\end{mdframed}

\section{Grammar}
\subsection{More on \itb{stare} and \itb{essere}}
\begin{itemize}[nosep]
    \item \itb{stare} to be, to stay: \iti{io sto}, \iti{tu stai}. 
	    Remember, \iti{stare} is never followed by an adjective.
    \item \itb{essere} to be: \iti{io sono}, \iti{iu sei}
    \item \itb{avere} to be: \iti{io ho}, \iti{tu hai}
\end{itemize}

\paragraph{Examples}
\begin{itemize}
    \item I am happy because this evening I am eating spaghetti carbonara. \\
    \iti{Io sono felice perché questa sera io sto mangiando spaghetti carbonara.}
    \item You are not happy because you are cooking spaghetti carbonara. \\
    \iti{Tu non sei felice perché tu stai cucinando spaghetti carbonara.}
    \item I am buying a coffee this morning. \\
    \iti{Io sto comprando un caffè questa mattina.}
    \item You are tired, because you cook a lot. \\
    \iti{Tu sei stanca, perché tu stai cucinando molto.}
    \item I am not happy because you have a lot of pasta and I do not. \\
    \iti{Io non sono felice perché tu hai molta pasta e io non ho.}
\end{itemize}

\subsection{Possessive adjectives}
Possessive adjectives `my', `his' or `hers', `your', etc. are \iti{mio}, \iti{suo}, 
and \iti{tuo} in their Italian base forms. They change ending just like other adjectives 
to match the noun being discussed. The phrase, ``The my car is a beautiful car'' is 
incorrect in English, but would be correct in Italian. Most of the time, we use articles 
with possessive adjectives. 
\begin{mdframed}
    \textbf{Exception} When we talk about family members, we omit the article.
\end{mdframed}

\paragraph{Examples}
\begin{itemize}
    \item My car is a beautiful car. \\
    \iti{La mia macchina è una bella macchina.}
    \item My brother is tired. \\
    \iti{Mio fratello è stanco.}
    \item I am happy because I am driving my new car. \\
    \iti{Io sono felice, perché io sto guidando la mia nuova macchina.}
    \item I am working with my computer this morning. \\
    \iti{Io sto lavorando con il mio computer questa mattina.}
\end{itemize}

\section{Vocabulary}
\subsection{Words}
\begin{table}[ht] %[ht] % place table roughly [h]ere or else at the [t]op of page
\centering
    \begin{tabular}{ r l  | r l  | r l  }
          car     & \iti{macchina} & to cook & \iti{cucinare}   & to buy    & \iti{comprare} \\
          morning & \iti{mattina}  & is      & \iti{è}          & and       & \iti{e} \\
          ceiling & \iti{soffitto} & roof    & \iti{tetto}      & wall      & \iti{muro} \\
          gelato  & \iti{gelato}   & because & \iti{perché}     & drawer    & \iti{cassetto} \\
          road    & \iti{strada}   & kitchen & \iti{cucina}     & class     & \iti{classe} \\
          husband & \iti{marito}   & wife    & \iti{moglie}     & delizioso & \iti{delicious} \\
	  small   & \iti{piccolo}  & angry   & \iti{arrabbiato} & name      & \iti{nome}
    \end{tabular}
    \caption{Words to practice (Jan 31)}
    \label{tab:tVocabJan31}
\end{table}

\section{Homework}
\begin{enumerate}
    \item I'm happy because I'm studying Italian. \\
    \iti{Io sono felice perché io sto studiando italiano.}
    \item I'm not feeling well because I don't have a new car. \\
    \iti{Io non sto bene perché io non ho una nuova macchina.}
    \item I'm tired, because I'm drinking a lot of wine. \\
    \iti{Io sono stanco, perché io sto bevando molto vino.}
    \item I'm tired, because I'm working today. \\
    \iti{Io sono stanco, perché io sto lavorando oggi.}
    \item My brother isn't doing well, because my car is very small. \\
    \iti{Mio fratello non sta bene, perché la mia auto è molto piccola.}
    \item My wife is happy, because you are eating your pasta. \\
    \iti{Mia moglie è contenta, perché tu stai mangiando la tua pasta.}
    \item My sister is feeling unwell, because you are cooking spaghetti. \\
    \iti{Mia sorella sta male, perché tu stai cucinando spaghetti.}
    \item Aren't you driving your new car today? \\
    \iti{Tu non stai guidando la tua nuova auto oggi?}
    \item I'm buying my coffee. \\
    \iti{Io sto comprando il mio caffè.}
    \item I'm fine, because my husband is talking. \\
    \iti{Io sto bene, perché mio marito sta parlando.}
\end{enumerate}
%----------------------------------------------------------------------------------------------------
% END OF CHAPTER
%----------------------------------------------------------------------------------------------------
